\documentclass[12pt]{article}

\usepackage{sbc-template}
\usepackage{graphicx,url}
\usepackage[brazil]{babel}
\usepackage[utf8]{inputenc}

\sloppy

\title{Segmentação de Instância de Animais de Estimação com YOLOv8: um Estudo de Caso no Dataset Oxford-IIIT Pet}

\author{Katelyn Moura, Diego Rodrigues, Idelmar Júnior,\\ Renan Costa, Thiago Novaes}

\address{Curso Técnico/Superior em Informática -- Instituto Federal do Piauí (IFPI)\\
  Campus Picos -- Picos, PI -- Brasil
  \email{\{katelyn.moura,diego.rodrigues,idelmar.junior,renan.costa,thiago.novaes\}@aluno.ifpi.edu.br}
}

\begin{document}

\maketitle

\begin{resumo}
  A segmentação é uma das tarefas fundamentais da visão computacional, permitindo a
  classificação de imagens no nível de pixel. Este trabalho apresenta a aplicação da
  arquitetura YOLOv8-seg para a tarefa de segmentação de instância de animais de
  estimação (cães e gatos), utilizando o dataset público Oxford-IIIT Pet. Como o dataset
  original disponibiliza apenas máscaras densas (\textit{trimaps}), foi desenvolvido um
  pipeline de conversão para o formato de polígonos exigido pela arquitetura. O modelo foi
  treinado por \textit{transfer learning} a partir de pesos pré-treinados no COCO, e avaliado
  por métricas de \textit{mean Average Precision} (mAP) de máscara. Os resultados demonstram
  [preencher após o treinamento], evidenciando a viabilidade do uso de arquiteturas de
  detecção com cabeça de segmentação para problemas de granularidade fina com poucos
  recursos computacionais.
\end{resumo}

\begin{abstract}
  Segmentation is one of the fundamental tasks in computer vision, enabling pixel-level
  image classification. This work presents the application of the YOLOv8-seg architecture
  to the instance segmentation of pets (cats and dogs), using the public Oxford-IIIT Pet
  dataset. Since the original dataset only provides dense masks (trimaps), a conversion
  pipeline was developed to produce the polygon format required by the architecture. The
  model was trained via transfer learning from COCO-pretrained weights and evaluated using
  mask mean Average Precision (mAP) metrics. Results show [to be filled in after training],
  demonstrating the feasibility of using detection architectures with a segmentation head
  for fine-grained problems under limited computational resources.
\end{abstract}


\section{Introdução}

A visão computacional compreende um conjunto de tarefas voltadas à extração de informação
semântica a partir de imagens digitais, das quais três são consideradas fundamentais:
classificação, detecção de objetos e segmentação. Enquanto a classificação indica \textit{o
que} está presente em uma imagem e a detecção localiza objetos por meio de caixas
delimitadoras, a segmentação é a tarefa mais granular, atribuindo uma classe a \textbf{cada
pixel} da imagem \cite{ronneberger2015unet}. Dentro da segmentação, distingue-se ainda a
segmentação semântica, que rotula pixels por classe sem diferenciar objetos individuais, da
segmentação de instância, que além de classificar também separa cada ocorrência de um
objeto com sua própria máscara.

Aplicações de segmentação são numerosas: em imagens médicas, permite delimitar órgãos e
lesões; em veículos autônomos, permite distinguir via, calçada e pedestres; em edição de
imagem, permite o recorte automático de objetos do fundo. Este trabalho tem como objetivo
aplicar uma arquitetura de segmentação de instância a um problema de domínio simples e bem
documentado -- a segmentação de animais de estimação em fotografias -- de modo a explorar,
na prática, todo o pipeline de um projeto de \textit{deep learning}: preparação de dados,
treinamento por \textit{transfer learning}, avaliação quantitativa e análise qualitativa dos
resultados.

O restante deste artigo está organizado da seguinte forma: a Seção~\ref{sec:teoria} apresenta
a fundamentação teórica sobre segmentação e a arquitetura utilizada; a
Seção~\ref{sec:metodologia} descreve o dataset, o pipeline de preparação dos dados e os
detalhes de treinamento; a Seção~\ref{sec:resultados} apresenta e discute os resultados
obtidos; e a Seção~\ref{sec:conclusao} conclui o trabalho e aponta direções futuras.


\section{Fundamentação Teórica}
\label{sec:teoria}

\subsection{Segmentação semântica e de instância}

Redes convolucionais do tipo \textit{encoder-decoder}, como a U-Net \cite{ronneberger2015unet},
consolidaram-se como abordagem clássica para segmentação semântica: um caminho de
contração (\textit{encoder}) extrai características em múltiplas escalas, enquanto um
caminho de expansão (\textit{decoder}) reconstrói a resolução espacial original, auxiliado
por conexões de salto (\textit{skip connections}) que preservam detalhes finos. Essa
abordagem, no entanto, não distingue instâncias individuais de uma mesma classe quando há
múltiplos objetos sobrepostos ou próximos na imagem.

\subsection{A família YOLO e o YOLOv8-seg}

A família YOLO (\textit{You Only Look Once}) \cite{redmon2016yolo} popularizou a detecção de
objetos em estágio único (\textit{single-stage}), predizendo simultaneamente, em uma única
passagem pela rede, as caixas delimitadoras, as classes e as confianças de todos os objetos
de uma imagem -- em contraste com abordagens de dois estágios, mais precisas porém mais
lentas. Versões mais recentes da arquitetura, mantidas pela Ultralytics
\cite{ultralytics2023yolov8}, estendem esse princípio à segmentação de instância
(\textit{YOLOv8-seg}): além da caixa e da classe, a rede prediz coeficientes que combinam
protótipos de máscara compartilhados, gerando uma máscara de pixels específica para cada
instância detectada, mantendo a eficiência computacional característica da família.


\section{Metodologia}
\label{sec:metodologia}

\subsection{Dataset}

Foi utilizado o \textit{Oxford-IIIT Pet Dataset} \cite{parkhi2012cats}, composto por 7.349
imagens de 37 raças de cães e gatos, cada uma acompanhada de uma máscara de segmentação em
formato \textit{trimap}, na qual cada pixel é rotulado como pertencente ao animal (pet), ao
fundo, ou à borda de contorno entre ambos. O dataset é distribuído publicamente e possui
divisão oficial entre um conjunto de treino/validação e um conjunto de teste.

\subsection{Conversão de máscaras para polígonos}

Como o YOLOv8-seg espera anotações no formato de polígono normalizado (e não máscaras
densas), foi implementado um pipeline de conversão: para cada imagem, os pixels rotulados
como pet ou borda no trimap foram unidos em uma máscara binária de primeiro plano; em
seguida, o contorno externo dessa máscara foi extraído com o algoritmo de
Suzuki e Abe \cite{suzuki1985contours}, implementado na biblioteca OpenCV
(\texttt{cv2.findContours}), e normalizado pelas dimensões da imagem para gerar o rótulo no
formato exigido pela arquitetura, com uma única classe (\textit{pet}). O conjunto de
treino/validação original foi então subdividido em treino (85\%) e validação (15\%),
mantendo o conjunto de teste oficial intacto para a avaliação final.

\subsection{Arquitetura e treinamento}

Foi utilizado o modelo \textit{YOLOv8n-seg}, a variante mais leve da família, com pesos
pré-treinados no conjunto de dados COCO \cite{lin2014coco}, adaptados ao problema por
\textit{transfer learning}. O treinamento foi conduzido com [\textit{X}] épocas, tamanho de
\textit{batch} de 16 imagens e resolução de entrada de $640\times640$ pixels, utilizando os
hiperparâmetros padrão da biblioteca \texttt{ultralytics} (otimizador SGD/Adam com
\textit{learning rate} adaptativo e função de perda combinada de caixa, classe e máscara). O
treinamento foi executado em ambiente Google Colab, com GPU gratuita (NVIDIA T4).

\subsection{Métricas de avaliação}

O desempenho do modelo foi avaliado no conjunto de teste por meio de métricas padrão de
detecção/segmentação: \textit{mean Average Precision} com limiar de sobreposição de máscara
de 0,5 (mAP50) e a média entre os limiares de 0,5 a 0,95 (mAP50-95), além de precisão e
\textit{recall}.


\section{Resultados e Discussão}
\label{sec:resultados}

A Tabela~\ref{tab:resultados} apresenta as métricas obtidas no conjunto de teste após o
treinamento do modelo.

\begin{table}[h]
\centering
\caption{Resultados no conjunto de teste}
\label{tab:resultados}
\begin{tabular}{lc}
\hline
\textbf{Métrica} & \textbf{Valor} \\
\hline
mAP50 (máscara)     & [preencher] \\
mAP50-95 (máscara)  & [preencher] \\
Precisão            & [preencher] \\
Recall               & [preencher] \\
\hline
\end{tabular}
\end{table}

% Figura de exemplo: adicionar após gerar as amostras com `src/predict.py`
% \begin{figure}[h]
%   \centering
%   \includegraphics[width=.8\textwidth]{figuras/exemplo_predicao.png}
%   \caption{Exemplo de predição: imagem original, máscara prevista e sobreposição.}
%   \label{fig:exemplo}
% \end{figure}

[Discussão a preencher: comparar os resultados quantitativos com o esperado, comentar
qualitativamente os acertos e erros observados nas Figuras~\ref{fig:exemplo}, e discutir as
principais dificuldades do desenvolvimento (dados, hardware, hiperparâmetros).]

Uma limitação observada é que o Oxford-IIIT Pet contém, em sua grande maioria, uma única
instância de animal por imagem, o que reduz a vantagem prática da segmentação de instância
sobre uma abordagem puramente semântica -- ponto relevante para trabalhos futuros que
explorem datasets com múltiplos animais por cena.


\section{Conclusão}
\label{sec:conclusao}

Este trabalho apresentou a aplicação da arquitetura YOLOv8-seg à tarefa de segmentação de
instância de animais de estimação, cobrindo todo o pipeline de um projeto de \textit{deep
learning}: desde a conversão de anotações do formato original do dataset para o formato
exigido pela arquitetura, passando pelo treinamento por \textit{transfer learning}, até a
avaliação quantitativa e qualitativa dos resultados. [Concluir com uma síntese dos principais
achados após a análise dos resultados.]

Como trabalhos futuros, sugere-se: (i) a comparação com arquiteturas de segmentação semântica
pura, como a U-Net, sob as mesmas condições experimentais; (ii) a avaliação em datasets com
múltiplas instâncias por imagem, de modo a explorar melhor a capacidade de segmentação de
instância do YOLOv8-seg; e (iii) a análise do custo-benefício entre variantes do modelo
(\textit{nano}, \textit{small}, \textit{medium}) em termos de acurácia \textit{versus} tempo
de inferência.


\bibliographystyle{sbc}
\bibliography{referencias}

\end{document}
